\chapter{Referencial teórico}
\label{cap:fundamentacao}

Este capítulo estabelece a base conceitual do trabalho em quatro movimentos. A
Seção~\ref{sec:pdn} descreve o comportamento eletromagnético da rede de
distribuição de energia e deriva as relações analíticas que serão posteriormente
utilizadas como restrições físicas. A Seção~\ref{sec:sparam} formaliza a
representação por parâmetros de espalhamento e sua conversão para impedância. A
Seção~\ref{sec:surrogate} revisa o estado da arte em modelos substitutos
aplicados a integridade de sinal e potência. A Seção~\ref{sec:pinn} apresenta o
paradigma de aprendizado informado por física e delimita a lacuna que este
trabalho pretende ocupar.

\section{A rede de distribuição de energia como estrutura eletromagnética}
\label{sec:pdn}

\subsection{Função e critério de projeto}

A PDN tem por função entregar corrente aos circuitos integrados mantendo a tensão
de alimentação dentro de uma janela de tolerância, tipicamente de $\pm 5\%$ em
torno do valor nominal. Como a corrente demandada por um circuito digital varia
abruptamente a cada transição lógica, a especificação de projeto é formulada no
domínio da frequência por meio de uma \emph{impedância-alvo}
\citep{swaminathan2020}:

\begin{equation}
Z_{\text{target}} = \frac{V_{dd} \cdot \delta}{I_{\max}},
\label{eq:ztarget}
\end{equation}

\noindent onde $V_{dd}$ é a tensão de alimentação, $\delta$ é a tolerância
relativa admitida e $I_{\max}$ é a máxima variação de corrente. O critério de
projeto exige que $|Z(f)| < Z_{\text{target}}$ ao longo de toda a faixa de
frequências em que o circuito integrado apresenta conteúdo espectral relevante.
Para um processador moderno, esse critério pode significar manter a impedância
abaixo de alguns miliohms de corrente contínua até centenas de megahertz — uma
exigência severa, cuja violação em qualquer ponto da faixa compromete o
desempenho.

\subsection{O par de planos como cavidade ressonante}

Em placas multicamadas, a PDN é implementada por pares de planos de cobre
separados por um dielétrico de espessura $h$, muito menor que as dimensões
laterais $a$ e $b$. Essa geometria constitui um guia de ondas de placas paralelas
cujas paredes laterais, por serem interfaces com o ar, comportam-se
aproximadamente como paredes magnéticas. Sob a hipótese $h \ll \lambda$, apenas
modos $\text{TM}_{mn0}$ propagam-se, e a estrutura admite solução analítica
fechada. As frequências de ressonância são \citep{swaminathan2020}:

\begin{equation}
f_{mn} = \frac{c_0}{2\sqrt{\varepsilon_r}}
         \sqrt{\left(\frac{m}{a}\right)^2 + \left(\frac{n}{b}\right)^2},
\qquad m,n \in \mathbb{Z}_{\geq 0},\; (m,n) \neq (0,0),
\label{eq:modos}
\end{equation}

\noindent onde $c_0$ é a velocidade da luz no vácuo e $\varepsilon_r$ a
permissividade relativa do dielétrico. Em uma cavidade pouco amortecida, a
impedância vista entre os planos apresenta máximos pronunciados nessas
frequências. É precisamente esse comportamento que estabelece a ponte entre
integridade de potência e compatibilidade eletromagnética: nos picos de
ressonância a cavidade acumula energia e a irradia pelas bordas abertas da placa,
transformando o par de planos em uma antena de fenda distribuída. Convém
registrar, desde já, que a observabilidade desses máximos em uma grandeza
específica — a autoimpedância de uma porta, por exemplo — depende do
amortecimento introduzido pelas vias e do acoplamento entre a porta e a
distribuição modal de campo. A Seção~\ref{sec:preliminar} examina essa questão
sobre os dados efetivamente utilizados.

No extremo oposto da faixa, quando $f \rightarrow 0$, o comprimento de onda
excede largamente as dimensões da placa e a cavidade comporta-se como um
capacitor concentrado de valor

\begin{equation}
C_{\text{plate}} = \varepsilon_0 \varepsilon_r \frac{a\,b}{h},
\label{eq:cplate}
\end{equation}

\noindent de modo que $|Z(f)| \approx 1/(2\pi f\, C_{\text{plate}})$, uma reta de
inclinação $-20$\,dB/década em escala logarítmica.

As Equações~\ref{eq:modos} e~\ref{eq:cplate} constituem o repertório físico
disponível em forma fechada para este problema. Ambas dependem exclusivamente de
parâmetros que integram o vetor de entrada do modelo substituto, o que as torna
candidatas naturais a restrições de treinamento. Qual delas é de fato utilizável
depende, contudo, de sua observabilidade na grandeza predita — questão empírica
que a Seção~\ref{sec:preliminar} resolve sobre os dados. O comportamento
intermediário da resposta, governado pelo acoplamento entre modos, pelas perdas
no dielétrico e no cobre e, sobretudo, pelo efeito do arranjo de vias que perfura
os planos, não admite forma fechada e constitui, justamente, o que se pretende
aprender dos dados.

\subsection{O papel do arranjo de vias}

Vias que atravessam os planos perturbam a cavidade de duas maneiras. Vias de
referência conectadas ao plano de terra atuam como curtos-circuitos distribuídos,
deslocando e amortecendo os modos ressonantes. Vias de sinal, por sua vez,
atravessam os planos através de aberturas circulares (\textit{antipads}) que
constituem descontinuidades de impedância e caminhos de acoplamento entre
cavidades adjacentes. O efeito líquido depende de forma não trivial do raio da
via, do raio do \textit{antipad}, do passo do arranjo e de sua posição relativa
às distribuições modais de campo — dependência que motiva o recurso a modelos
empíricos \citep{hillebrecht2024}.

\section{Representação por parâmetros de espalhamento}
\label{sec:sparam}

Em frequências de micro-ondas, a caracterização de estruturas com múltiplos
acessos é feita por parâmetros de espalhamento, que relacionam ondas de tensão
incidentes e refletidas em cada porta. Para uma estrutura de $P$ portas, a matriz
$\mathbf{S}(f) \in \mathbb{C}^{P \times P}$ descreve completamente o
comportamento linear do sistema. A conversão para a matriz de impedância é dada
por

\begin{equation}
\mathbf{Z}(f) = \sqrt{\mathbf{Z}_0}\,
                \left(\mathbf{I} + \mathbf{S}(f)\right)
                \left(\mathbf{I} - \mathbf{S}(f)\right)^{-1}
                \sqrt{\mathbf{Z}_0},
\label{eq:s2z}
\end{equation}

\noindent onde $\mathbf{Z}_0$ é a matriz diagonal de impedâncias de referência
das portas, tipicamente $50\,\Omega$. A autoimpedância $Z_{11}(f)$, obtida do
elemento $(1,1)$ dessa matriz, corresponde à impedância que o circuito integrado
enxerga ao demandar corrente da PDN, e é a grandeza de interesse primário deste
trabalho.

Três propriedades da resposta são fisicamente obrigatórias e serão usadas como
restrições no Capítulo~\ref{cap:metodologia}. A \emph{passividade} exige que a
estrutura não gere energia, o que se traduz na condição de que todos os valores
singulares de $\mathbf{S}(f)$ sejam menores ou iguais à unidade, ou
equivalentemente que $\operatorname{Re}\{Z_{11}(f)\} \geq 0$ para toda
frequência. A \emph{causalidade} exige que a resposta não anteceda a excitação, o
que vincula as partes real e imaginária da resposta pelas relações de
Kramers-Kronig, implementáveis numericamente pela transformada de Hilbert. A
\emph{reciprocidade}, válida para meios isotrópicos e sem materiais girotrópicos,
exige que $\mathbf{S}$ seja simétrica. \citet{torun2020} demonstram que
redes neurais treinadas apenas para minimizar erro numérico violam rotineiramente
as duas primeiras propriedades, gerando modelos inúteis para simulação
subsequente no domínio do tempo.

\section{Modelos substitutos em integridade de sinal e potência}
\label{sec:surrogate}

Um modelo substituto é uma aproximação computacionalmente barata de uma função
custosa de avaliar. No contexto eletromagnético, substitui o solver por um
regressor treinado sobre um conjunto de simulações previamente executadas
\citep{swaminathan2020}.

\subsection{Panorama de abordagens}

O trabalho de referência de \citet{swaminathan2020} sistematiza a aplicação de
aprendizado de máquina a problemas de integridade de sinal e potência em
encapsulamento, cobrindo desde redes neurais rasas até processos gaussianos, e
identifica a escassez de dados como o gargalo dominante. \citet{zhang2021}
demonstram a viabilidade em escala ao treinar uma rede profunda sobre mais de um
milhão de placas geradas sinteticamente por método de elementos de contorno,
obtendo predição de impedância de PDN em 0,1\,s — cerca de cem vezes mais rápido
que o método de contorno e aproximadamente cinco mil vezes mais rápido que
simulação de onda completa. Esse resultado estabelece o teto de desempenho
alcançável quando dados são abundantes, e simultaneamente delimita a condição em
que a abordagem funciona: um orçamento de simulação inacessível à maioria dos
grupos de pesquisa.

Estudos comparativos recentes indicam que a escolha do regressor depende
criticamente do regime de dados. Processos gaussianos com núcleo anisotrópico
apresentam desempenho superior quando as amostras são escassas, ao passo que
redes neurais os superam largamente em conjuntos volumosos. Essa dependência
justifica a estratégia adotada neste trabalho, que estabelece uma bateria de
referências abrangendo ambos os regimes antes de introduzir o modelo informado
por física.

\subsection{Otimização apoiada por modelos substitutos}

Uma vez disponível um substituto rápido, torna-se viável embuti-lo em laços de
otimização. A linha mais explorada na literatura é o posicionamento de
capacitores de desacoplamento por aprendizado por reforço profundo
\citep{zhang2020decap,duan2024}, em que um agente aprende a selecionar posições e
valores de capacitância que satisfaçam a impedância-alvo com o menor número de
componentes. Essas abordagens compartilham uma premissa restritiva: tratam a
geometria da placa como fixa e otimizam apenas os componentes adicionados a
posteriori. O presente trabalho adota a perspectiva complementar, otimizando os
parâmetros do próprio empilhamento — decisão que antecede o posicionamento de
capacitores e que, por ocorrer mais cedo no fluxo de projeto, tem custo marginal
de implementação muito inferior.

Para a etapa de otimização, adotam-se algoritmos evolutivos multiobjetivo, campo
em que se concentra parte substancial da produção do grupo de pesquisa no qual o
trabalho se insere \citep{coelho2010,takara2026}. A escolha por metaheurísticas
populacionais, em detrimento de métodos baseados em gradiente, decorre da
natureza mista do espaço de busca — que combina variáveis contínuas, como
espessura de dielétrico e permissividade, e variáveis discretas, como número de
camadas e contagem de vias — e da presença de múltiplos ótimos locais associados
à estrutura modal da cavidade.

\subsection{Bases de dados públicas}

A ausência de dados padronizados dificultou por muito tempo a comparação entre
métodos na área. A SI/PI-Database, mantida pelo \textit{Institut für Theoretische
Elektrotechnik} da Universidade Técnica de Hamburgo, foi criada especificamente
para suprir essa lacuna \citep{schierholz2021}. Ela reúne vinte e cinco
estruturas distintas de PCB, cada uma acompanhada de um arquivo tabular de
parâmetros de projeto e dos correspondentes arquivos Touchstone com a resposta
eletromagnética. \citet{hillebrecht2024} ampliaram substancialmente a base com
três subconjuntos de aproximadamente quinze mil variações de arranjos de vias
simulados até 60\,GHz, e demonstraram sua aplicação à análise de dependências
entre parâmetros de PCB e a frequência de relógio de PCI Gen 6. A base é
gerada por modelagem baseada em física e validada por comparação com solver de
onda completa, o que confere aos rótulos qualidade suficiente para servirem de
referência de treinamento.

\section{Aprendizado informado por física}
\label{sec:pinn}

\subsection{Formulação}

O paradigma de redes neurais informadas por física, consolidado por
\citet{raissi2019}, propõe incorporar conhecimento físico diretamente à função de
perda, transformando o treinamento em um problema de otimização sujeito a
restrições brandas. A perda assume a forma composta

\begin{equation}
\mathcal{L} = \underbrace{\mathcal{L}_{\text{dados}}}_{\text{ajuste às amostras}}
            + \sum_{k} \lambda_k \,
              \underbrace{\mathcal{L}_{\text{física},k}}_{\text{violação de restrição}},
\label{eq:pinn}
\end{equation}

\noindent em que cada termo $\mathcal{L}_{\text{física},k}$ penaliza o desvio em
relação a uma lei conhecida e $\lambda_k$ pondera sua importância relativa. O
efeito prático é uma redução do espaço de hipóteses: soluções que ajustam bem os
dados de treinamento mas violam a física tornam-se inacessíveis ao otimizador. O
benefício é máximo, portanto, exatamente no regime em que os dados são escassos —
situação em que a evidência empírica, isoladamente, não basta para descartar tais
soluções.

\subsection{Aplicação a integridade de sinal e potência}

A adoção do paradigma em integridade de sinal e potência ainda é incipiente e
concentra-se no pós-processamento. \citet{torun2020} propõem camadas
dedicadas de imposição de causalidade e passividade, aplicadas à saída da rede,
que garantem consistência física das respostas preditas mediante as relações de
Kramers-Kronig e as propriedades espectrais dos parâmetros de espalhamento.
Trabalhos mais recentes exploram modelos substitutos guiados por física para
interconexões verticais, combinando restrições físicas com otimização
multiobjetivo por enxame de partículas.

Permanece, entretanto, um espaço pouco explorado. A literatura em integridade de
potência não tira proveito de uma circunstância particularmente favorável do
problema da cavidade retangular: a existência de soluções analíticas exatas nos
dois extremos da faixa de análise, dadas pelas Equações~\ref{eq:modos}
e~\ref{eq:cplate}. Diferentemente de aplicações em que a restrição física é uma
equação diferencial cuja avaliação demanda diferenciação automática custosa,
aqui as restrições são expressões algébricas de custo desprezível, calculáveis
diretamente a partir do vetor de entrada. Essa característica torna o problema um
candidato especialmente adequado ao paradigma informado por física, e é sobre
essa observação que se assenta a contribuição pretendida por este trabalho.

\section{Síntese e posicionamento}
\label{sec:sintese}

O Quadro~\ref{quad:posicionamento} resume o posicionamento da proposta em relação
aos trabalhos discutidos.

\begin{table}[htb]
\centering
\caption{Posicionamento da proposta em relação ao estado da arte}
\label{quad:posicionamento}
\footnotesize
\begin{tabular}{@{}p{0.20\linewidth}p{0.16\linewidth}p{0.17\linewidth}p{0.17\linewidth}p{0.20\linewidth}@{}}
\toprule
\textbf{Trabalho} & \textbf{Alvo} & \textbf{Regime de dados} & \textbf{Física
incorporada} & \textbf{Otimização} \\
\midrule
\citet{zhang2021} & $Z(f)$ de PDN & $>10^6$ amostras & não & não \\[0.3em]
\citet{torun2020} & $\mathbf{S}(f)$ & moderado & pós-processamento &
não \\[0.3em]
\citet{zhang2020decap}; \citet{duan2024} & posição de capacitores & simulação
em laço & não & aprendizado por reforço \\[0.3em]
\citet{hillebrecht2024} & análise de dependências & $\approx 15\times 10^3$
amostras & modelagem dos dados & não \\[0.3em]
\textbf{Esta proposta} & $Z(f)$ de PDN & $985$ amostras & na função de perda &
evolutiva multiobjetivo \\
\bottomrule
\end{tabular}
\vspace{0.3em}
\\{\footnotesize FONTE: O autor (2026).}
\end{table}

A lacuna que o trabalho pretende ocupar pode ser enunciada de forma direta: os
métodos de alto desempenho existentes dependem de orçamentos de simulação
inacessíveis, enquanto os métodos que incorporam física o fazem apenas na saída
do modelo. Este trabalho investiga se a incorporação de física \emph{durante} o
treinamento permite operar com acurácia útil em um regime de dados três ordens de
grandeza menor.
